\section{Related Work}
\label{sec:related}

\paragraph{Keeping documents true to programs.} The aspiration behind
machine-checked courseware descends from literate programming
\citep{knuth1984literate}, which interleaves exposition with the executable
source, and from executable-documentation practice such as \texttt{doctest},
which runs the examples embedded in prose. BabyWhale extends this line in two
directions: the checked unit is not an embedded example but a \emph{claim
about a maintained external system} (a displayed excerpt, an exercise's
semantics, a measured number), and the checking is adversarial---our audit
(\S\ref{sec:eval}) shows the errors that matter are semantic and survive
example-level testing.

\paragraph{Educational systems.}

Table~\ref{tab:comparison} situates BabyWhale among widely used educational
resources. nanoGPT and Raschka's book teach GPT-style pre-training (the latter
adding fine-tuning) with clean, runnable code, but stop before the modern architecture and the
post-training/serving stages. The Annotated Transformer
and labml.ai pioneered side-by-side prose-and-code presentation, which we adopt---but their code is
re-implemented \emph{for} the article, so nothing prevents divergence from any
maintained system. Spinning Up established the
theory$\to$pseudocode$\to$minimal-code bridge for deep RL; our
``theory-to-code'' tables generalize this to systems mechanisms (a page table,
a batching protocol) where the left column is a rule rather than an equation.
TinyTorch \citep{tinytorch2026} is closest in pedagogy: students build a
framework under an autograder with real-task milestones. BabyWhale inherits
its build-and-verify stance but inverts the scope: rather than rebuilding
tensors and autograd, learners build the \emph{LLM layer}---MLA, routing, DPO,
speculative acceptance---on top of MLX, and are graded against a maintained,
tested production-shaped stack rather than a scaffold.

Two properties are, to our knowledge, unique to BabyWhale: (i) full-lifecycle
coverage including mid-training, verifiable-reward RL, and a
continuous-batching server, on one shared codebase; and (ii) CI-enforced
consistency between the teaching material and the implementation
(\S\ref{sec:course}).

\begin{table*}[t]
\centering
\small
\begin{tabular}{@{}lcccccc@{}}
\toprule
 & \textbf{nanoGPT} & \textbf{Raschka} & \textbf{Annot.\,Tr.\,/\,labml} & \textbf{Spinning Up} & \textbf{TinyTorch} & \textbf{BabyWhale} \\
\midrule
Pre-training & \cmark & \cmark & -- & -- & \cmark & \cmark \\
Mid-training (context/anneal) & -- & -- & -- & -- & -- & \cmark \\
SFT + preference opt.\ (DPO) & -- & \cmark & -- & -- & -- & \cmark \\
RL w/ verifiable rewards & -- & -- & -- & \cmark\,(RL) & -- & \cmark \\
Modern arch.\ (MLA/MoE/MTP) & -- & -- & partial & -- & -- & \cmark \\
Serving (paged KV, batching, spec.) & -- & -- & -- & -- & KV cache & \cmark \\
Quantization & -- & -- & -- & -- & \cmark & \cmark \\
Graded build exercises & -- & -- & -- & exercises & \cmark & \cmark\,(23) \\
Graded vs.\ \emph{maintained} system & -- & -- & -- & -- & -- & \cmark \\
Runnable measured payoffs & -- & -- & -- & -- & \cmark & \cmark \\
CI drift guard on docs & -- & -- & -- & -- & -- & \cmark \\
\bottomrule
\end{tabular}
\caption{Educational LLM/ML systems compared. ``\cmark'' indicates first-class,
documented support. BabyWhale is the only system covering the full lifecycle
on one codebase with machine-checked teaching material.}
\label{tab:comparison}
\end{table*}
